\part{2D Ising Model}

\section{Thermodynamic limit}

\begin{exercise}[Boundary terms, subadditivity, and the existence of the pressure]
Let $d\ge1$.  A box is a finite set $\Lambda\subset\Z^d$; $\Lambda_L=\{1,\dots,L\}^d$;
For state limits put $\Lambda_L^{\rm c}=\{-L,\dots,L\}^d$. $E(\Lambda)$ is the set of nearest-neighbour edges with both ends in
$\Lambda$; $\partial\Lambda$ is the set of sites of $\Lambda$ with a
nearest neighbour outside $\Lambda$; $\abs\Lambda$ is the number of sites.
Let $\Omega=\{\pm1\}^{\Z^d}$ with the product $\sigma$-algebra.  For
$\eta\in\Omega$ the Ising Hamiltonian in $\Lambda$ with boundary condition
$\eta$ is
\[
 H^\eta_\Lambda(\sigma)=-J\sum_{ij\in E(\Lambda)}\sigma_i\sigma_j-h\sum_{i\in\Lambda}\sigma_i
 -J\sum_{\substack{i\in\Lambda,\ j\notin\Lambda\\ i\sim j}}\sigma_i\eta_j,
 \qquad\sigma\in\Omega_\Lambda=\{\pm1\}^\Lambda,
\]
with partition function $Z^\eta_\Lambda(\beta,h)$ and Gibbs measure
$\mu^\eta_\Lambda$ (the section Gibbs measures of Part~I); the free
boundary condition $\eta=\emptyset$ omits the last sum, giving
$Z^\emptyset_\Lambda$.  Put $K=\beta J$.
\begin{enumerate}[label=(\alph*)]
\item Prove $\abs{\log Z^\eta_\Lambda-\log Z^\emptyset_\Lambda}\le 2dK\abs{\partial\Lambda}$
for every $\eta$, and, if $\Lambda=\Lambda_1\sqcup\Lambda_2$,
\[
 e^{-K\abs{E(\Lambda_1,\Lambda_2)}}Z^\emptyset_{\Lambda_1}Z^\emptyset_{\Lambda_2}
 \le Z^\emptyset_\Lambda\le
 e^{K\abs{E(\Lambda_1,\Lambda_2)}}Z^\emptyset_{\Lambda_1}Z^\emptyset_{\Lambda_2},
\]
where $E(\Lambda_1,\Lambda_2)$ is the set of edges between $\Lambda_1$ and $\Lambda_2$.
\item Prove that $\psi_k=2^{-kd}\log Z^\emptyset_{\Lambda_{2^k}}$ satisfies
$\abs{\psi_{k+1}-\psi_k}\le dK2^{-k}$, hence converges to a limit
$\psi(\beta,h)$, and that
$\abs{\Lambda_L}^{-1}\log Z^\emptyset_{\Lambda_L}\to\psi(\beta,h)$ as $L\to\infty$
(tile $\Lambda_L$ by translates of $\Lambda_{2^k}$ and a remainder of at most
$d\,2^kL^{d-1}$ sites, and use $\abs{\log Z^\emptyset_\Lambda}\le(dK+\beta\abs h+\log2)\abs\Lambda$).
\item A sequence of boxes $\Lambda^{(n)}$ is a van Hove sequence if
$\abs{\Lambda^{(n)}}\to\infty$ and, for every $r$,
\[
 \frac{\abs{\{i\in\Lambda^{(n)}:\operatorname{dist}(i,\Z^d\setminus\Lambda^{(n)})\le r\}}}{\abs{\Lambda^{(n)}}}\longrightarrow0 .
\]  Prove that $\abs{\Lambda^{(n)}}^{-1}\log Z^{\eta_n}_{\Lambda^{(n)}}\to\psi(\beta,h)$
for every van Hove sequence and every choice of boundary conditions $\eta_n$
(including $\emptyset$), and for the torus $\T_L=(\Z/L\Z)^d$ with periodic
boundary conditions.  The limit $\psi(\beta,h)$ is the pressure; the free
energy density is $f=-\beta^{-1}\psi$.
\item Prove $dK\le\psi(\beta,h)-\beta\abs h\le dK+\log2$, that $\psi$ is even in
$h$, and that $\psi(\beta,h)-\beta\abs h\to dK$ as $\abs h\to\infty$.
\end{enumerate}
\end{exercise}

\begin{exercise}[Convexity of the pressure, magnetization, and one-sided derivatives]
Keep the notation of the previous exercise and regard $\psi$ as a function of
$(\beta,\beta h)$ as in the section $\log Z$, cumulants, entropy, free energy
of Part~I.
\begin{enumerate}[label=(\alph*)]
\item Prove that $\psi$ is convex in $(\beta,\beta h)\in(0,\infty)\times\R$
and continuous, and that for fixed $\beta$ the function $\beta h\mapsto\psi$ is
convex, even, and nondecreasing on $[0,\infty)$.
\item Prove that the one-sided derivatives
$m^\pm(\beta,h)=\partial^\pm\psi/\partial(\beta h)$ exist everywhere, are
nondecreasing in $h$, satisfy $m^-\le m^+$, and coincide except at countably
many $h$; the magnetization $m(\beta,h)=\partial\psi/\partial(\beta h)$ is
defined where they coincide, and $\abs{m}\le1$.
\item Prove that if $\psi$ is differentiable in $\beta h$ at $(\beta,h)$ then
for every van Hove sequence and every boundary condition
\[
 \frac1{\abs{\Lambda^{(n)}}}\Bigl\langle\sum_{i\in\Lambda^{(n)}}\sigma_i\Bigr\rangle^{\eta_n}_{\Lambda^{(n)}}
 \longrightarrow m(\beta,h)
\]
(a sequence of convex functions converging pointwise has derivatives
converging at every point of differentiability of the limit).
\item Define the spontaneous magnetization
$m^*(\beta)=\lim_{h\downarrow0}m(\beta,h)$.  Prove that the limit exists, that
$m^*(\beta)=m^+(\beta,0)=-m^-(\beta,0)\ge0$, and that $\psi$ is differentiable
in $h$ at $h=0$ if and only if $m^*(\beta)=0$.
\item Prove that $\beta\mapsto\beta f(\beta,0)=-\psi(\beta,0)$ is concave,
that its one-sided derivatives $u^\pm(\beta)$ (the internal energy per site)
exist and are nonincreasing, and that for $\beta_1<\beta_2$,
$\psi(\beta_2,0)-\psi(\beta_1,0)=-\int_{\beta_1}^{\beta_2}u(\beta)\dd\beta$
with $u=u^+$ almost everywhere.
\end{enumerate}
\end{exercise}

\section{Boundary conditions}

\begin{exercise}[The FKG lattice condition, the four functions theorem, Holley's inequality, and the FKG inequality]
Let $\Lambda$ be a finite set and order $\{\pm1\}^\Lambda$ pointwise, with
$(\sigma\vee\sigma')_i=\max(\sigma_i,\sigma'_i)$ and
$(\sigma\wedge\sigma')_i=\min(\sigma_i,\sigma'_i)$.  A function
$f:\{\pm1\}^\Lambda\to\R$ is increasing if $\sigma\le\sigma'$ implies
$f(\sigma)\le f(\sigma')$; an event is increasing if its indicator is.  A
strictly positive measure $\mu$ on $\{\pm1\}^\Lambda$ satisfies the FKG
lattice condition if
\[
 \mu(\sigma\vee\sigma')\,\mu(\sigma\wedge\sigma')\ge\mu(\sigma)\,\mu(\sigma')
 \qquad\text{for all }\sigma,\sigma' .
\]
Two positive measures $\mu_1,\mu_2$ satisfy the Holley condition if
$\mu_1(\sigma\vee\sigma')\mu_2(\sigma\wedge\sigma')\ge\mu_1(\sigma)\mu_2(\sigma')$
for all $\sigma,\sigma'$.  For probability measures, $\mu_2\preceq\mu_1$
(stochastic domination) means $\E_{\mu_2}f\le\E_{\mu_1}f$ for every increasing $f$.
\begin{enumerate}[label=(\alph*)]
\item (Four functions, one site.)  Let $\alpha,\beta,\gamma,\delta:\{\pm1\}\to[0,\infty)$
satisfy $\alpha(s)\beta(s')\le\gamma(s\vee s')\delta(s\wedge s')$ for all
$s,s'$.  Prove
$\bigl(\alpha(+)+\alpha(-)\bigr)\bigl(\beta(+)+\beta(-)\bigr)\le\bigl(\gamma(+)+\gamma(-)\bigr)\bigl(\delta(+)+\delta(-)\bigr)$.
(If $\gamma(+)\delta(-)=0$ argue directly; otherwise normalize
$\gamma(+)=\delta(-)=1$, replace $\gamma(-)$ and $\delta(+)$ by the smaller
values $\alpha(-)\beta(-)$ and $\alpha(+)\beta(+)$, and reduce to
$(1-\alpha(-)\beta(+))(1-\alpha(+)\beta(-))\ge0$.)
\item (Four functions theorem.)  Let $\alpha,\beta,\gamma,\delta:\{\pm1\}^\Lambda\to[0,\infty)$
satisfy $\alpha(\sigma)\beta(\sigma')\le\gamma(\sigma\vee\sigma')\delta(\sigma\wedge\sigma')$
for all $\sigma,\sigma'$.  Prove by induction on $\abs\Lambda$ that
$\alpha(\{\pm1\}^\Lambda)\beta(\{\pm1\}^\Lambda)\le\gamma(\{\pm1\}^\Lambda)\delta(\{\pm1\}^\Lambda)$,
where $\alpha(A)=\sum_{\sigma\in A}\alpha(\sigma)$: fix a site $i$, define
$\alpha'(\tau)=\alpha(\tau,+)+\alpha(\tau,-)$ on $\{\pm1\}^{\Lambda\setminus\{i\}}$
and similarly $\beta',\gamma',\delta'$, and verify the hypothesis for the
primed functions using (a).  Deduce
$\alpha(A)\beta(B)\le\gamma(A\vee B)\delta(A\wedge B)$ for all
$A,B\subset\{\pm1\}^\Lambda$, where $A\vee B=\{\sigma\vee\sigma':\sigma\in A,\sigma'\in B\}$.
\item (Holley's inequality.)  If probability measures $\mu_1,\mu_2$ satisfy
the Holley condition, prove $\mu_2\preceq\mu_1$ (apply (b) to
$\alpha=f\mu_2$, $\beta=\mu_1$, $\gamma=f\mu_1$, $\delta=\mu_2$ for
$f\ge0$ increasing).
\item (FKG inequality.)  If a probability measure $\mu$ satisfies the lattice
condition, prove $\E_\mu(fg)\ge\E_\mu f\,\E_\mu g$ for all increasing $f,g$,
and $\E_\mu(fg)\le\E_\mu f\,\E_\mu g$ if $f$ is increasing and $g$ decreasing.
\item Prove that if $\mu$ satisfies the lattice condition, then so does every
conditional measure $\mu(\,\cdot\mid\sigma_i=s)$ on $\{\pm1\}^{\Lambda\setminus\{i\}}$,
and that the pair $\mu(\,\cdot\mid\sigma_i=+1)$, $\mu(\,\cdot\mid\sigma_i=-1)$
satisfies the Holley condition.
\end{enumerate}
\end{exercise}

\begin{exercise}[Monotone couplings from the Holley condition]
Let $\mu_1,\mu_2$ be probability measures on $\{\pm1\}^\Lambda$ satisfying
the Holley condition.  Prove, by induction on $\abs\Lambda$, that there is a
probability measure $P$ on $\{\pm1\}^\Lambda\times\{\pm1\}^\Lambda$ with
marginals $\mu_1$ and $\mu_2$ and $P(\sigma^{(1)}\ge\sigma^{(2)})=1$: choose
a site $i$; from Holley's inequality $\mu_1(\sigma_i=+1)\ge\mu_2(\sigma_i=+1)$,
so the two Bernoulli laws of $\sigma_i$ admit a monotone coupling; and for
$s\ge s'$ the conditional measures $\mu_1(\,\cdot\mid\sigma_i=s)$ and
$\mu_2(\,\cdot\mid\sigma_i=s')$ again satisfy the Holley condition (compare
the points $(\sigma,s)$ and $(\sigma',s')$).  Deduce that for any
$f:\{\pm1\}^\Lambda\to\R$,
\[
 \abs{\E_{\mu_1}f-\E_{\mu_2}f}\le\sum_{i\in\Lambda}\operatorname{osc}_i(f)\,
 \frac{\E_{\mu_1}\sigma_i-\E_{\mu_2}\sigma_i}2,
\]
where $\operatorname{osc}_i(f)$ is the maximal change of $f$ under a flip of
$\sigma_i$.
\end{exercise}

\begin{exercise}[The Ising model satisfies the lattice condition; Griffiths inequalities via duplicate variables]
Let $\Lambda\subset\Z^d$ be a box.  The boundary conditions $+$, $-$ are
$\eta\equiv+1$, $\eta\equiv-1$; $\emptyset$ is free; periodic means the
Hamiltonian of the torus.  More generally admit real boundary values
$\eta_j\in\R$ and site-dependent fields $h_i\in\R$ in $H^\eta_\Lambda$.
\begin{enumerate}[label=(\alph*)]
\item Prove that for $a,a',b,b'\in\{\pm1\}$,
$(a\vee a')(b\vee b')+(a\wedge a')(b\wedge b')\ge ab+a'b'$, and deduce that
$\mu^\eta_\Lambda$ satisfies the lattice condition for every $J\ge0$, every
$\eta\in\R^{\Z^d\setminus\Lambda}$, and all fields $h_i\in\R$; likewise the
free and periodic measures.
\item Prove that if $\eta\le\eta'$ pointwise and $h_i\le h'_i$, then the pair
$\mu^{\eta'}_\Lambda(h')$, $\mu^\eta_\Lambda(h)$ satisfies the Holley
condition, hence $\mu^\eta_\Lambda(h)\preceq\mu^{\eta'}_\Lambda(h')$; in
particular $\mu^-_\Lambda\preceq\mu^\eta_\Lambda\preceq\mu^+_\Lambda$ for every
$\eta\in\Omega$, and $\mu^\emptyset_\Lambda\preceq\mu^+_\Lambda$.
\item (Griffiths I.)  For $A\subset\Lambda$ put $\sigma_A=\prod_{i\in A}\sigma_i$.
Expanding every factor $e^{K\sigma_i\sigma_j}$, $e^{\beta h_i\sigma_i}$, and
$e^{K\sigma_i\eta_j}$ in its power series, prove
$\langle\sigma_A\rangle^\eta_\Lambda\ge0$ whenever $J\ge0$, $h_i\ge0$, and
$\eta_j\ge0$ for all $i,j$.
\item (Griffiths II, Ginibre's argument.)  Under the same hypotheses and for
$A,B\subset\Lambda$, write
$\langle\sigma_A\sigma_B\rangle-\langle\sigma_A\rangle\langle\sigma_B\rangle
=(Z^\eta_\Lambda)^{-2}\sum_{\sigma,\sigma'}\sigma_A(\sigma_B-\sigma'_B)
e^{-\beta H^\eta_\Lambda(\sigma)-\beta H^\eta_\Lambda(\sigma')}$
and substitute $\sigma_i=(x_i+y_i)/2$, $\sigma'_i=(x_i-y_i)/2$.  Prove that
$\sigma_A(\sigma_B-\sigma'_B)$, $\sigma_i\sigma_j+\sigma'_i\sigma'_j=(x_ix_j+y_iy_j)/2$,
and $\sigma_i+\sigma'_i=x_i$ are polynomials in $(x,y)$ with nonnegative
coefficients, that the exponential of a polynomial with nonnegative
coefficients has a power series with nonnegative coefficients, and that
$\sum_{\sigma_i,\sigma'_i\in\{\pm1\}}x_i^ay_i^b\ge0$ for all $a,b\ge0$.
Conclude
\[
 \langle\sigma_A\sigma_B\rangle^\eta_\Lambda\ge\langle\sigma_A\rangle^\eta_\Lambda\langle\sigma_B\rangle^\eta_\Lambda ,
\]
and that $\langle\sigma_A\rangle^\eta_\Lambda$ is nondecreasing in each
coupling $K$, in each field $\beta h_i\ge0$, and in each boundary value
$\eta_j\ge0$.
\end{enumerate}
\end{exercise}

\begin{exercise}[Probability measures on spin products and conditional expectations]
For the countable product $\Omega=\{\pm1\}^{\Z^d}$ a cylinder
set fixes finitely many spins. Suppose probability laws on all
finite spin sets are consistent under restriction.
\begin{enumerate}[label=(\alph*)]
\item Define their common finitely additive law on the algebra
of cylinder sets. Use compactness to show that a decreasing
sequence of cylinder sets with empty intersection is
eventually empty. Extend the resulting premeasure and prove
uniqueness on the product sigma-algebra by approximation
with cylinder functions.
\item For a probability measure $\mu$, the space
$L^2(\mu)$ is the completion of bounded measurable functions
modulo zero-norm functions for
$\|f\|_2^2=\int|f|^2\dd\mu$.
For a sub-sigma-algebra $\mathcal A$, define
$\E_\mu(f\mid\mathcal A)$ as the orthogonal projection
onto its closed subspace of $\mathcal A$-measurable functions.
Prove the integral characterization against bounded
$\mathcal A$-measurable functions, positivity, and the tower
property; extend to $L^1$ by truncation.
\item Prove that orthogonal projections onto a decreasing
sequence of closed Hilbert subspaces converge in norm to
the projection onto their intersection, by comparing the
squared norms and applying the Cauchy criterion.
Deduce the corresponding convergence of bounded conditional
expectations; treat sigma-algebras modulo null sets.
\end{enumerate}
\end{exercise}

\begin{exercise}[Monotonicity in the volume and the infinite-volume limits \texorpdfstring{$\mu^\pm$}{mu+-}]
Let $\Lambda\subset\Lambda'$ be boxes in $\Z^d$.  A function
$f:\Omega\to\R$ is local if it depends on finitely many coordinates.
\begin{enumerate}[label=(\alph*)]
\item Prove $\mu^+_{\Lambda'}\preceq\mu^+_\Lambda$ and
$\mu^-_\Lambda\preceq\mu^-_{\Lambda'}$ (condition $\mu^+_{\Lambda'}$ on
$\sigma|_{\Lambda'\setminus\Lambda}$ and use the finite-volume consistency of
Part~I together with the previous exercise).
\item Prove that $\langle\sigma_0\rangle^+_\Lambda$ is nondecreasing in $h$
and, for $h\ge0$, nondecreasing in $\beta$; and that
$\langle\sigma_0\rangle^+_\Lambda\ge\langle\sigma_0\rangle^\emptyset_\Lambda\ge0$
for $h\ge0$.
\item Prove that every local $f$ is a difference of two increasing local
functions.  Deduce that for every local $f$ the limits
$\mu^+(f)=\lim_{L\to\infty}\mu^+_{\Lambda_L^{\rm c}}(f)$ and $\mu^-(f)=\lim\mu^-_{\Lambda_L^{\rm c}}(f)$
exist, are the same along every sequence of boxes increasing to $\Z^d$, and
define (Kolmogorov extension) probability measures $\mu^\pm$ on $\Omega$ with
$\mu^-\preceq\mu^+$; that $\mu^\pm$ are invariant under translations of $\Z^d$
and under the lattice symmetries fixing the origin; and that
$\mu^+(f)=\mu^-(f\circ(-1))$ at $h=0$, where $(-1)$ is global spin reversal.
\item Prove that $\mu^\pm$ satisfy the FKG inequality for increasing local
functions, and that $\langle\sigma_A\rangle^+\ge0$,
$\langle\sigma_A\sigma_B\rangle^+\ge\langle\sigma_A\rangle^+\langle\sigma_B\rangle^+$
for $h\ge0$.
\end{enumerate}
\end{exercise}

\section{Gibbs states}

\begin{exercise}[The Gibbs specification, DLR equations, and the Gibbs property of \texorpdfstring{$\mu^\pm$}{mu+-}]
Fix $\beta>0$, $h\in\R$.  For a box $\Lambda$ let $\mathcal F_{\Lambda^c}$ be
the $\sigma$-algebra on $\Omega$ generated by $(\sigma_j)_{j\notin\Lambda}$.
The Gibbs specification is the family of probability kernels
$\gamma_\Lambda(A\mid\eta)=\sum_{\tau\in\Omega_\Lambda}\mu^\eta_\Lambda(\tau)\,
\mathbf 1_A(\tau\eta)$, $A$ measurable, $\eta\in\Omega$, where
$\tau\eta$ agrees with $\tau$ on $\Lambda$ and with $\eta$ off $\Lambda$;
for bounded measurable $f$ put
$(\gamma_\Lambda f)(\eta)=\sum_\tau\mu^\eta_\Lambda(\tau)f(\tau\eta)$.  A
probability measure $\mu$ on $\Omega$ is a Gibbs state (for $\beta,h$) if it
satisfies the DLR equations
\[
 \mu(A\mid\mathcal F_{\Lambda^c})=\gamma_\Lambda(A\mid\cdot\,)\quad\mu\text{-a.s.},
 \qquad\text{for every box }\Lambda\text{ and measurable }A ,
\]
equivalently $\mu(\gamma_\Lambda f)=\mu(f)$ for all bounded measurable $f$ and
all $\Lambda$.  Write $\mathcal G(\beta,h)$ for the set of Gibbs states.
\begin{enumerate}[label=(\alph*)]
\item Prove that $\gamma_\Lambda f$ is $\mathcal F_{\Lambda^c}$-measurable,
that $\gamma_\Lambda(gf)=g\,\gamma_\Lambda f$ for $g$
$\mathcal F_{\Lambda^c}$-measurable, that $\gamma_\Lambda f$ is local
whenever $f$ is local, and that $\gamma_{\Lambda'}\gamma_\Lambda=\gamma_{\Lambda'}$
for $\Lambda\subset\Lambda'$ (the finite-volume consistency of Part~I).
Prove the equivalence of the two formulations of the DLR equations.
\item Prove that $\mu^+$ and $\mu^-$ are Gibbs states: for local $f$ and
$\Lambda\subset\Lambda_L^{\rm c}$, $\mu^+_{\Lambda_L^{\rm c}}(\gamma_\Lambda f)=\mu^+_{\Lambda_L^{\rm c}}(f)$,
and $\gamma_\Lambda f$ is local.
\item Prove that $\mathcal G(\beta,h)$ is convex and closed in the weak
topology of probability measures on the compact space $\Omega$, hence
compact, and that every Gibbs state $\mu$ satisfies
$\mu^-\preceq\mu\preceq\mu^+$ on increasing local functions
(write $\mu(f)=\mu(\gamma_\Lambda f)$ and bound $\gamma_\Lambda f$ pointwise
by $\mu^\pm_\Lambda(f)$ using the section Boundary conditions).
\item Prove that if $\mu$ is a translation-invariant Gibbs state then
$\mu(\sigma_0)\in[\langle\sigma_0\rangle^-,\langle\sigma_0\rangle^+]$, where
$\langle\sigma_0\rangle^\pm=\mu^\pm(\sigma_0)$.
\end{enumerate}
\end{exercise}

\begin{exercise}[Extremal Gibbs states and tail triviality]
The tail $\sigma$-algebra is $\mathcal T=\bigcap_{L}\mathcal F_{(\Lambda_L^{\rm c})^c}$.
A probability measure $\mu$ is tail trivial if $\mu(A)\in\{0,1\}$ for all
$A\in\mathcal T$.  A Gibbs state $\mu$ is extremal if $\mu=t\nu_1+(1-t)\nu_2$
with $\nu_1,\nu_2\in\mathcal G(\beta,h)$ and $t\in(0,1)$ implies $\nu_1=\nu_2=\mu$.
\begin{enumerate}[label=(\alph*)]
\item Let $\mu\in\mathcal G(\beta,h)$ and let $\nu\in\mathcal G(\beta,h)$ be
absolutely continuous with respect to $\mu$, with density $g$.  For a box
$\Lambda$ let $g_\Lambda=\E_\mu(g\mid\mathcal F_{\Lambda^c})$.  Prove
$\nu(A)=\mu(g\,\gamma_\Lambda(A))=\mu(g_\Lambda\gamma_\Lambda(A))=\mu(g_\Lambda\mathbf 1_A)$
for every measurable $A$, and conclude $g=g_\Lambda$ $\mu$-a.s. for every box $\Lambda$.
\item Deduce that for every $c\in\Q$ the event $\{g>c\}$ coincides
$\mu$-a.s. with the tail event $\limsup_L\{g_{\Lambda_L^{\rm c}}>c\}$.  Prove: if
$\mu$ is tail trivial then $g$ is $\mu$-a.s. constant, hence $\nu=\mu$.
\item Prove that a tail-trivial Gibbs state is extremal (if
$\mu=t\nu_1+(1-t)\nu_2$ then $\nu_1\ll\mu$).  Conversely, if
$\mu\in\mathcal G(\beta,h)$ and $A\in\mathcal T$ with $0<\mu(A)<1$, prove that
$\mu(\,\cdot\mid A)$ and $\mu(\,\cdot\mid A^c)$ are Gibbs states (use
$\gamma_\Lambda(f\mathbf 1_A)=\mathbf 1_A\gamma_\Lambda f$), so that $\mu$ is not
extremal.  Conclude that a Gibbs state is extremal if and only if it is tail
trivial.
\item Prove that a tail-trivial probability measure $\mu$ on $\Omega$ has
the mixing property
\[
 \lim_{L\to\infty}\sup_{B\in\mathcal F_{(\Lambda_L^{\rm c})^c}}
 \abs{\mu(A\cap B)-\mu(A)\mu(B)}=0
\]
for every measurable $A$ (approximate $\mathbf 1_A$ by local functions and use
that $\E_\mu(\mathbf 1_A\mid\mathcal F_{(\Lambda_L^{\rm c})^c})$ converges in $L^1(\mu)$
to $\E_\mu(\mathbf 1_A\mid\mathcal T)=\mu(A)$; prove the convergence first for bounded variables by the decreasing-subspace projection argument, and then use $L^1$ approximation).
\end{enumerate}
\end{exercise}

\begin{exercise}[Uniqueness of the Gibbs state and the magnetization criterion]
Fix $\beta>0$ and $h\in\R$.
\begin{enumerate}[label=(\alph*)]
\item Prove that the following are equivalent: (i) $\mathcal G(\beta,h)$ has
exactly one element; (ii) $\mu^+=\mu^-$; (iii)
$\langle\sigma_0\rangle^+=\langle\sigma_0\rangle^-$.  For (iii)$\Rightarrow$(ii)
apply the coupling bound of the section Boundary conditions to
$\mu^+_\Lambda,\mu^-_\Lambda$, then let $\Lambda\uparrow\Z^d$ and use the
translation invariance of $\mu^\pm$.
\item Prove that at $h=0$ the condition (iii) reads $\langle\sigma_0\rangle^+_{\beta,0}=0$.
\item Prove that $\beta\mapsto\langle\sigma_0\rangle^+_{\beta,0}$ is
nondecreasing, so that
$\beta_c=\inf\{\beta>0:\langle\sigma_0\rangle^+_{\beta,0}>0\}\in[0,\infty]$
satisfies: $\mathcal G(\beta,0)$ is a singleton for $\beta<\beta_c$ and has
at least the two elements $\mu^+\ne\mu^-$ for $\beta>\beta_c$.
\item Prove that $\mu^+$ and $\mu^-$ are extremal Gibbs states, in the
following two steps.  First, if $\mu^+=t\nu_1+(1-t)\nu_2$ with
$\nu_1,\nu_2\in\mathcal G$, then $\nu_1\preceq\mu^+$ and $\nu_2\preceq\mu^+$ by
the previous exercise, and $t\nu_1(f)+(1-t)\nu_2(f)=\mu^+(f)$ for increasing
$f$ forces $\nu_1(f)=\nu_2(f)=\mu^+(f)$; second, increasing local functions
determine a measure.
\end{enumerate}
\end{exercise}

\section{Correlation functions}

\begin{exercise}[Correlation functions, truncated correlations, and their monotonicity]
For a Gibbs state $\mu$ and a finite $A\subset\Z^d$ the $n$-point function
is $\langle\sigma_A\rangle=\mu(\sigma_A)$, $n=\abs A$; the truncated two-point
function is $\langle\sigma_x;\sigma_y\rangle=\langle\sigma_x\sigma_y\rangle-\langle\sigma_x\rangle\langle\sigma_y\rangle$.
\begin{enumerate}[label=(\alph*)]
\item Prove that for $h\ge0$, $\langle\sigma_A\rangle^+$ is nondecreasing in
$\beta$ and in $h$, that $\langle\sigma_x;\sigma_y\rangle^+\ge0$, and that
$\langle\sigma_x;\sigma_y\rangle^+\le1-(\langle\sigma_0\rangle^+)^2$.
\item Prove that at $h=0$, $\langle\sigma_A\rangle^\emptyset_\Lambda=0$ for
$\abs A$ odd, $\langle\sigma_A\rangle^+\ge\langle\sigma_A\rangle^\emptyset_\Lambda\ge0$
for $A\subset\Lambda$, and $\langle\sigma_x\sigma_y\rangle^+\ge(\langle\sigma_0\rangle^+)^2$.
\item Prove that $\mu^+$ is mixing in the sense of the section Gibbs states,
so that $\langle\sigma_0;\sigma_x\rangle^+\to0$ as $\abs x\to\infty$, and hence
$\lim_{\abs x\to\infty}\langle\sigma_0\sigma_x\rangle^+=(\langle\sigma_0\rangle^+)^2$.
\item Define the correlation length $\xi(\beta,h)\in[0,\infty]$ by
\[
 \frac1{\xi}=\liminf_{n\to\infty}\Bigl(-\frac1n\log\langle\sigma_0;\sigma_{ne_1}\rangle^+\Bigr),
\]
with $-\log0=+\infty$ in the limit inferior.  Prove that $\xi$ is unchanged when $e_1$ is replaced by any other unit
vector of $\Z^d$, that $\xi\le1/c$ whenever
$\langle\sigma_0;\sigma_x\rangle^+\le Ce^{-c\abs x}$ for all $x$.
\end{enumerate}
\end{exercise}

\begin{exercise}[Exponential decay at high temperature from the path expansion]
Let $\Lambda\subset\Z^d$ be a box, $h=0$, and $t=\tanh K$.  Adjoin a ghost
site $g$ joined by one edge to each boundary site $i\in\partial\Lambda$ for
each of its neighbours outside $\Lambda$, so that the $+$ boundary term is
$K\sum_{i}n_i\sigma_i=K\sum_{ig}\sigma_i\sigma_g$ with $\sigma_g=+1$ fixed.
Write $\bar E$ for the edge set of $\Lambda$ together with the ghost edges.
\begin{enumerate}[label=(\alph*)]
\item Prove $e^{K\sigma_i\sigma_j}=\cosh K\,(1+t\sigma_i\sigma_j)$, and by
expanding the product over edges and summing over $\sigma\in\Omega_\Lambda$,
\[
 Z^+_\Lambda=2^{\abs\Lambda}(\cosh K)^{\abs{\bar E}}\sum_{E'\subset\bar E:\ \deg_{E'}(i)\text{ even }\forall i\in\Lambda}t^{\abs{E'}},
 \qquad
 \langle\sigma_0\rangle^+_\Lambda=\frac{\sum_{E':\ \deg_{E'}(i)\text{ odd iff }i=0}t^{\abs{E'}}}
 {\sum_{E':\ \deg_{E'}(i)\text{ even }\forall i\in\Lambda}t^{\abs{E'}}},
\]
where degrees are counted at sites of $\Lambda$ only; the same formulas with
$\bar E$ replaced by $E(\Lambda)$ give $Z^\emptyset_\Lambda$ and
$\langle\sigma_0\sigma_x\rangle^\emptyset_\Lambda$ (odd degree exactly at $0$ and $x$).
\item Prove that every $E'$ in the numerator for $\langle\sigma_0\rangle^+_\Lambda$
contains a self-avoiding path from $0$ to $g$ (the total number of odd-degree
vertices of a finite graph is even), and every $E'$ in the numerator of
$\langle\sigma_0\sigma_x\rangle^\emptyset_\Lambda$ contains a self-avoiding path from
$0$ to $x$.  Removing a canonically chosen such path $\gamma$ from $E'$ leaves
a set with all degrees even; deduce
\[
 \langle\sigma_0\rangle^+_\Lambda\le\sum_{\gamma:0\to g}t^{\abs\gamma},\qquad
 \langle\sigma_0\sigma_x\rangle^\emptyset_\Lambda\le\sum_{\gamma:0\to x}t^{\abs\gamma},
\]
the sums running over self-avoiding paths.
\item Prove that the number of paths of length $n$ from $0$ in $\Z^d$ is at
most $(2d)^n$, and conclude that for $2d\tanh K<1$,
\[
 \langle\sigma_0\rangle^+_\Lambda\le\frac{(2dt)^{\operatorname{dist}(0,\partial\Lambda)+1}}{1-2dt},\qquad
 \langle\sigma_0\sigma_x\rangle^\emptyset_\Lambda\le\frac{(2dt)^{\abs x_1}}{1-2dt},
\]
hence $\langle\sigma_0\rangle^+_{\beta,0}=0$, the Gibbs state at $(\beta,0)$ is
unique, and $\beta_c\ge\beta_0$ with $\tanh(\beta_0J)=1/(2d)$.
\item Prove that for $2d\tanh K<1$ the unique Gibbs state satisfies
$\langle\sigma_0\sigma_x\rangle\le(2dt)^{\abs x_1}/(1-2dt)$ (every weak limit
point of the free measures $\mu^\emptyset_{\Lambda_L^{\rm c}}$ is a Gibbs state), so that
$\xi(\beta,0)\le1/\log(1/(2d\tanh K))$.
\end{enumerate}
\end{exercise}

\begin{exercise}[Two-point functions on the cylinder from the spectrum of the transfer matrix]
Let $T_L$ be the transfer matrix of the section Transfer matrices of Part~I
on the row space $\mathcal H_L$, with spectral decomposition
$T_L=\sum_a\lambda_aP_a$, $\lambda_{\max}=\lambda_0>\lambda_1\ge\lambda_2\ge\dots>0$
($P_a$ are mutually orthogonal spectral projections, refined
by spin reversal when $h=0$, with their full multiplicities), and
let $\Sigma_j$ be multiplication by $\tau_j$ on $\mathcal H_L$.  Consider the
torus $\T_{N,L}$ and the sites $(0,0)$ and $(n,0)$ in rows $0$ and $n$.
\begin{enumerate}[label=(\alph*)]
\item Prove
\[
 \langle\sigma_{(0,0)}\sigma_{(n,0)}\rangle_{\T_{N,L}}
 =\frac{\sum_{a,b}\lambda_a^{n}\lambda_b^{N-n}\Tr(P_a\Sigma_0P_b\Sigma_0)}{\sum_a\lambda_a^N\Tr P_a},
 \qquad 0\le n\le N .
\]
\item Prove that as $N\to\infty$ with $n$ and $L$ fixed (the infinite
cylinder of circumference $L$),
\[
 \langle\sigma_{(0,0)}\sigma_{(n,0)}\rangle_{L}
 =\sum_{a}\Bigl(\frac{\lambda_a}{\lambda_{\max}}\Bigr)^n\Tr(P_0\Sigma_0P_a\Sigma_0),
 \]
\[
 \Tr(P_0\Sigma_0P_a\Sigma_0)\ge0,\qquad
 \sum_a\Tr(P_0\Sigma_0P_a\Sigma_0)=1 .
\]
\item At $h=0$ prove that $\Sigma_0$ anticommutes with the spin reversal $F$,
that $P_0$ lies in the $F=+1$ eigenspace, hence $\Tr(P_0\Sigma_0)=0$ and only
$a$ with $P_a$ in the $F=-1$ eigenspace contribute; conclude that
$\langle\sigma_{(0,0)};\sigma_{(n,0)}\rangle_L$ decays as
$e^{-n/\xi_L}$ with $1/\xi_L=\log(\lambda_{\max}/\lambda^{(-)}_{\max})$, where
$\lambda^{(-)}_{\max}$ is the largest eigenvalue $\lambda_a$ in the $F=-1$
sector with $\Tr(P_0\Sigma_0P_a\Sigma_0)\ne0$.
\end{enumerate}
\end{exercise}

\section{Phase transitions}

\begin{exercise}[Phase transitions: nonanalyticity of the pressure, nonuniqueness of Gibbs states, and their equivalence at zero field]
A phase transition occurs at $(\beta,h)$ in the thermodynamic sense if
$\psi$ is not real analytic at $(\beta,h)$, and in the sense of Gibbs states
if $\abs{\mathcal G(\beta,h)}>1$.
\begin{enumerate}[label=(\alph*)]
\item Prove that $h\mapsto\langle\sigma_0\rangle^+_{\beta,h}$ is
nondecreasing and right-continuous (a decreasing limit of continuous
nondecreasing functions of $h$ is upper semicontinuous).
\item Prove, for $0\le h<h'$ and every box $\Lambda$,
\[
\begin{aligned}
 \beta\int_h^{h'}\langle\sigma_0\rangle^+_{\beta,s}\dd s
 &\le\frac{\log Z^+_\Lambda(\beta,h')-\log Z^+_\Lambda(\beta,h)}{\abs\Lambda}\\
 &\le\beta\int_h^{h'}\Bigl(\langle\sigma_0\rangle^+_{\Lambda'_\ell,s}
 +\frac{\abs{\{i\in\Lambda:\operatorname{dist}(i,\Lambda^c)\le\ell\}}}{\abs\Lambda}\Bigr)\dd s .
\end{aligned}
\]
for every $\ell$, using $\langle\sigma_i\rangle^+_\Lambda\ge\langle\sigma_i\rangle^+$
and $\langle\sigma_i\rangle^+_\Lambda\le\langle\sigma_i\rangle^+_{i+\Lambda_\ell'}$
for the box $i+\Lambda'_\ell$ of side $2\ell+1$ centred at $i$ when it lies in
$\Lambda$.  Deduce $\psi(\beta,h')-\psi(\beta,h)=\beta\int_h^{h'}\langle\sigma_0\rangle^+_{\beta,s}\dd s$
and hence
\[
 m^*(\beta)=m^+(\beta,0)=\langle\sigma_0\rangle^+_{\beta,0},\qquad
 m(\beta,h)=\langle\sigma_0\rangle^+_{\beta,h}\ \text{ at every $h>0$ where $\psi$ is differentiable}.
\]
\item Conclude that at $h=0$ the following are equivalent: $\psi(\beta,\cdot)$
is differentiable at $0$; $m^*(\beta)=0$; $\mathcal G(\beta,0)$ is a
singleton.  Conclude also that $\beta_c$ of the section Gibbs states equals
$\sup\{\beta:m^*(\beta)=0\}$, that $m^*(\beta)>0$ for $\beta>\beta_c$, and that
for $d=1$, $\beta_c=\infty$ (the section Transfer matrices of Part~I).
\end{enumerate}
\end{exercise}

\begin{exercise}[The Peierls argument on \texorpdfstring{$\Z^2$}{Z2}]
Let $d=2$, $h=0$, and let $\Lambda\subset\Z^2$ be a box with $+$ boundary
condition.  The dual lattice is $(\Z^2)^*=\Z^2+(\tfrac12,\tfrac12)$; each edge
$ij$ of $\Z^2$ crosses exactly one dual edge $(ij)^*$.  For
$\sigma\in\Omega_\Lambda$ (extended by $+1$ outside $\Lambda$) let
$E^*(\sigma)$ be the set of dual edges $(ij)^*$ with $\sigma_i\ne\sigma_j$.  A
contour is a cycle (closed self-avoiding polygon) of the dual lattice; its
length $\abs\gamma$ is its number of edges; it surrounds $0$ if $0$ lies in
the bounded component of its complement in $\R^2$.
\begin{enumerate}[label=(\alph*)]
\item Prove that $E^*(\sigma)$ has even degree at every dual vertex.  Prove
that if $\sigma_0=-1$ then $E^*(\sigma)$ contains a contour surrounding $0$:
let $C$ be the connected component of $\{i:\sigma_i=-1\}$ containing $0$
(finite, contained in $\Lambda$), and take the boundary of the unbounded
component of the complement of the union of the closed unit squares centred
at the sites of $C$.
\item For a contour $\gamma\subset E^*(\sigma)$ let $\sigma^\gamma$ be
obtained by reversing every spin inside $\gamma$.  Prove that
$\sigma\mapsto\sigma^\gamma$ is an injection of $\{\sigma:\gamma\subset E^*(\sigma)\}$
into $\Omega_\Lambda$ with $H^+_\Lambda(\sigma^\gamma)=H^+_\Lambda(\sigma)-2J\abs\gamma$,
and deduce $\mu^+_\Lambda(\gamma\subset E^*(\sigma))\le e^{-2K\abs\gamma}$.
\item Prove that the number of contours of length $\ell$ surrounding $0$ is
at most $\tfrac\ell2\,3^{\ell-1}$ (such a contour crosses the half-line
$\{(x,0):x\ge0\}$ at a vertical dual edge within distance $\ell/2$ of $0$, and
a self-avoiding polygon is determined by its successive turns) and that
$\ell$ is even and $\ell\ge4$.
\item Conclude, with $y=9e^{-4K}<1$,
\[
 \mu^+_\Lambda(\sigma_0=-1)\le\sum_{\ell\ge4\text{ even}}\frac\ell2\,3^{\ell-1}e^{-2K\ell}
 =\frac13\cdot\frac{y^2(2-y)}{(1-y)^2},
 \qquad
 \langle\sigma_0\rangle^+_{\beta,0}\ge1-\frac23\cdot\frac{y^2(2-y)}{(1-y)^2},
\]
uniformly in $\Lambda$; hence $m^*(\beta)>0$ for $K\ge\tfrac34$ (where the
bound gives $m^*\ge0.318$), and $\beta_c(2)J\le\tfrac34$.  Combined with the
section Correlation functions, $\operatorname{artanh}(\tfrac14)\le\beta_cJ\le\tfrac34$
for $d=2$, so that $0<\beta_c<\infty$ and there is a phase transition at
$(\beta_c,0)$; and $\beta_c(d)<\infty$ for every $d\ge2$ (embed $\Z^2$ and use
monotonicity of $\langle\sigma_0\rangle^+$ in the couplings).
\end{enumerate}
\end{exercise}

\begin{exercise}[Critical exponents and the Rushbrooke inequality]
Suppose $0<\beta_c<\infty$ and $m^*(\beta_c)=0$.  Write $\tau=(\beta-\beta_c)/\beta_c$, and
regard $\psi$ as a function of $(\beta,\eta)$ with $\eta=\beta h$.  Where the
derivatives exist put $\chi=\beta\,\partial^2\psi/\partial\eta^2$
(the susceptibility $\partial m/\partial h$), $c=k_B\beta^2\,\partial^2\psi/\partial\beta^2$
(the specific heat per site; at $\eta=0$ the derivative at fixed $\eta$
coincides with the derivative at fixed $h=0$, so this agrees with Part~I).
The critical exponents are defined, when the limits exist, by
\begin{alignat*}{2}
 &c(\beta,0)\sim\abs\tau^{-\alpha}\ (\tau\to0),&\qquad
 &m^*(\beta)\sim\tau^{\beta_{\mathrm{mag}}}\ (\tau\downarrow0),\\
 &\chi(\beta,0)\sim\abs\tau^{-\gamma}\ (\tau\to0),&
 &m(\beta_c,h)\sim h^{1/\delta}\ (h\downarrow0),\\
 &\xi(\beta,0)\sim\abs\tau^{-\nu}\ (\tau\to0),&
 &\langle\sigma_0;\sigma_x\rangle_{\beta_c,0}\sim\abs x^{-(d-2+\eta_{\mathrm{corr}})}\ (\abs x\to\infty),
\end{alignat*}
in this display interpret the power laws as
$\log c/\log|\tau|\to-\alpha$ and the corresponding logarithmic
ratios for the other observables; retain the usual ratio-one
meaning of $\sim$ elsewhere. A logarithmic divergence of $c$ is
recorded as $\alpha=0$ (log).  Primed exponents $\alpha',\gamma'$ refer to
$\tau\downarrow0$.
\begin{enumerate}[label=(\alph*)]
\item Prove that the Hessian of $\psi$ in $(\beta,\eta)$ is positive
semidefinite wherever $\psi$ is $C^2$, hence
$c\,\chi\ge k_B\beta^3\bigl(\partial m/\partial\beta\bigr)^2$ there, where
$m=\partial\psi/\partial\eta$ and $\partial m/\partial\beta$ is taken at fixed $\eta$.
\item Assume: $\psi$ is $C^2$ on $\{\beta_c<\beta<\beta_c+\epsilon,\ 0<\eta<\epsilon\}$;
there are constants with $c\le A\tau^{-\alpha'}$ and
$\chi\le B\tau^{-\gamma'}$ on that region; $m(\beta,\eta)\to m^*(\beta)$ as
$\eta\downarrow0$ for each $\beta$; and $m^*(\beta)\ge D\tau^{\beta_{\mathrm{mag}}}$.
Prove, by integrating (a) in $\beta$ from $\beta_c$ to $\beta$ at fixed $\eta$ and
letting $\eta\downarrow0$, that
\[
 \alpha'+2\beta_{\mathrm{mag}}+\gamma'\ge2 .
\]
(If $\alpha'+\gamma'\ge2$ there is nothing to prove.)
\item For the Ising chain at $h=0$ prove
$\chi/(\beta\xi)\to2$ as $K\to\infty$.
Distinguish this zero-temperature relation from exponents
defined at a finite critical temperature.
\end{enumerate}
\end{exercise}

\section{Lee--Yang zeros}

\begin{exercise}[The partition function as a polynomial in the fugacity and the Asano contraction lemma]
Let $\Gamma=(V,E)$ be a finite graph, $K\ge0$, and for $\sigma\in\{\pm1\}^V$
let $n_-(\sigma)=\abs{\{i:\sigma_i=-1\}}$.  Put
\[
 P_\Gamma(z_1,\dots,z_n)=\sum_{\sigma}\Bigl(\prod_{ij\in E}e^{K\sigma_i\sigma_j}\Bigr)\prod_{i:\sigma_i=-1}z_i ,
 \qquad V=\{1,\dots,n\},
\]
a polynomial of degree at most one in each variable (multiaffine), and
$P_\Gamma(z)=P_\Gamma(z,\dots,z)$, so that
$Z_\Gamma(K,\beta h)=e^{\beta h\abs V}P_\Gamma(e^{-2\beta h})$ with $z=e^{-2\beta h}$.
Let $D=\{z\in\C:\abs z<1\}$.
\begin{enumerate}[label=(\alph*)]
\item Prove $P_\Gamma(z)=z^{\abs V}P_\Gamma(1/z)$ (spin reversal), that
$P_\Gamma$ has real nonnegative coefficients, leading and constant
coefficients $e^{K\abs E}$, and that the zero set of $P_\Gamma$ is invariant
under $z\mapsto1/z$ and under complex conjugation.
\item (Asano contraction lemma.)  Let $A(z_1,z_2)=a+bz_1+cz_2+dz_1z_2$ with
$a,b,c,d\in\C$ and $A(z_1,z_2)\ne0$ for all $(z_1,z_2)\in D\times D$.  Prove
that $a\ne0$; that $\abs{a+bz_1}\ge\abs{c+dz_1}$ for all $z_1\in\overline D$
and $\abs{a+cz_2}\ge\abs{b+dz_2}$ for all $z_2\in\overline D$; by
averaging the squares of these inequalities over $\abs{z_1}=1$ and
$\abs{z_2}=1$, that $\abs a\ge\abs d$; and conclude that $a+dz\ne0$ for
$z\in D$.
\item (Edge polynomials.)  Prove that
\[
 A_K(z,z')=e^{K}+e^{-K}z+e^{-K}z'+e^{K}zz'
 =\sum_{\sigma,\sigma'\in\{\pm1\}}e^{K\sigma\sigma'}
              z^{[\sigma=-1]}z'^{[\sigma'=-1]}
\]
has no zero in $D\times D$ for $K\ge0$: solving $A_K=0$ for $z'$, show that
$\abs{e^K+e^{-K}z}^2-\abs{e^{-K}+e^Kz}^2=2\sinh(2K)(1-\abs z^2)\ge0$.
\end{enumerate}
\end{exercise}

\begin{exercise}[The Lee--Yang circle theorem]
Keep the notation of the previous exercise.
\begin{enumerate}[label=(\alph*)]
\item Introduce a variable $z_{e,i}$ for each pair (edge $e$, endpoint $i$
of $e$), and put $Q=\prod_{e=ij\in E}A_K(z_{e,i},z_{e,j})$, a multiaffine
polynomial in the $2\abs E$ variables with no zero in the polydisc
$D^{2\abs E}$.  Define the contraction of two variables $z_{e,i},z_{e',i}$ at
the same site $i$ into one variable $z_i$ as the operation replacing
$a+bz_{e,i}+cz_{e',i}+dz_{e,i}z_{e',i}$ (with $a,b,c,d$ polynomials in the
remaining variables) by $a+dz_i$.  Prove, by (b) of the previous exercise
applied with the remaining variables fixed in the polydisc, that a
contraction of a multiaffine polynomial without zeros in the polydisc has no
zeros in the polydisc.
\item Prove that contracting, for each site $i$, all variables $z_{e,i}$
($e\ni i$) into $z_i$ turns $Q$ into $P_\Gamma(z_1,\dots,z_n)$ (a term of $Q$
survives all contractions exactly when, at every site, either all or none of
the variables $z_{e,i}$ occur, which is the consistency of the spin values
$\sigma_i$ across the edges at $i$); treat isolated sites separately.
Conclude that $P_\Gamma(z_1,\dots,z_n)\ne0$ on $D^n$.
\item (Circle theorem.)  Prove that all zeros of $P_\Gamma(z)$ lie on the
unit circle $\abs z=1$, for every finite graph $\Gamma$ and every $K\ge0$.
\item Prove that the same holds for the Ising model on a box $\Lambda$ with
free or periodic boundary conditions and with site-dependent couplings
$K_{ij}\ge0$, and that with $+$ boundary conditions the zeros of the
polynomial $e^{-\beta h\abs\Lambda}Z^+_\Lambda$ in $z=e^{-2\beta h}$ lie in
$\{\abs z\ge1\}$ (a boundary field $Kn_i\sigma_i$ rescales $z_i$ by $e^{-2Kn_i}$).
\end{enumerate}
\end{exercise}

\begin{exercise}[Consequences: analyticity of the pressure in the field and the density of zeros]
Let $\Lambda=\Lambda_L\subset\Z^d$ with free boundary conditions, and write
$P_\Lambda(z)=e^{K\abs{E(\Lambda)}}\prod_{k=1}^{\abs\Lambda}(z-e^{\ii\theta_k})$
with $\theta_k\in(-\pi,\pi]$ (the circle theorem).  Let
$\rho_\Lambda=\abs\Lambda^{-1}\sum_k\delta_{\theta_k}$, a probability measure on
the circle, symmetric under $\theta\mapsto-\theta$.
\begin{enumerate}[label=(\alph*)]
\item Prove that for $z\notin\{\abs z=1\}$,
\[
 \frac{\log Z_\Lambda(K,\beta h)}{\abs\Lambda}=\beta h+K\frac{\abs{E(\Lambda)}}{\abs\Lambda}
 +\int\log\abs{z-e^{\ii\theta}}\dd\rho_\Lambda(\theta),\qquad z=e^{-2\beta h}>0 ,
\]
and that every subsequential weak limit $\rho$ of $(\rho_{\Lambda_L})$ satisfies
$\psi(\beta,h)=\beta h+dK+\int\log\abs{z-e^{\ii\theta}}\dd\rho(\theta)$ for
all $h\ne0$ (the integrand is continuous in $\theta$ for $\abs z\ne1$).
\item Prove that $u(z)=\int\log\abs{z-e^{\ii\theta}}\dd\rho(\theta)$ is
harmonic on $D$ and on $\C\setminus\overline D$, and conclude that
$h\mapsto\psi(\beta,h)$ is real analytic on $(0,\infty)$ and on $(-\infty,0)$
for every $\beta>0$; hence in the thermodynamic sense a phase transition in
$h$ at fixed $\beta$ can occur only at $h=0$, and $m(\beta,h)$,
$\chi(\beta,h)$ are real analytic in $h\ne0$.
\item Prove that if there is $\delta>0$ with $\rho_{\Lambda_L}((-\delta,\delta))=0$
for all large $L$, then $\psi(\beta,\cdot)$ is real analytic in a
neighbourhood of $h=0$ and $m^*(\beta)=0$.
\item Suppose that a subsequential limit $\rho$ has, near $\theta=0$, a
density $g(\theta)$ with respect to $\dd\theta$ that is continuous at $0$.
Prove
\[
 m(\beta,h)=1-2z\,\partial_zu(z)=1-2z\int\frac{z-\cos\theta}{1-2z\cos\theta+z^2}\dd\rho(\theta),
 \qquad z=e^{-2\beta h},
\]
and, letting $z\uparrow1$ (Poisson kernel), $m^*(\beta)=2\pi g(0)$.
\end{enumerate}
\end{exercise}

\section{High/low-temperature expansions}

\begin{exercise}[High-temperature and low-temperature expansions and Kramers--Wannier duality on the torus]
Let $\T_L=(\Z/L\Z)^2$ with edge set $E$, $\abs E=2L^2$, $h=0$, $t=\tanh K$, and let
$K^*>0$ be defined by $\tanh K^*=e^{-2K}$, equivalently
$\sinh2K\sinh2K^*=1$.  The dual torus $\T_L^*$ has one vertex per face of
$\T_L$ and one edge $e^*$ crossing each edge $e$ of $\T_L$.
\begin{enumerate}[label=(\alph*)]
\item (High temperature.)  Prove
$Z_{\T_L}(K)=2^{L^2}(\cosh K)^{2L^2}\sum_{E'\subset E\text{ even}}t^{\abs{E'}}$,
the sum over edge sets with even degree at every vertex, and
$\langle\sigma_x\sigma_y\rangle=\sum_{E'\text{ odd exactly at }x,y}t^{\abs{E'}}/\sum_{E'\text{ even}}t^{\abs{E'}}$.
\item (Low temperature.)  For a configuration $\sigma$ on $\T_L$ let
$\Gamma(\sigma)\subset E^*$ be the set of dual edges crossing edges with
$\sigma_i\ne\sigma_j$.  Prove that $\Gamma(\sigma)$ is even, that
$\sigma\mapsto\Gamma(\sigma)$ is two-to-one onto its image, and that
$Z_{\T_L}(K)=2e^{2KL^2}\sum_{\Gamma\in\mathcal B}e^{-2K\abs\Gamma}$, where
$\mathcal B$ is the image: the even subsets of $E^*$ that are boundaries,
that is, those crossing every closed loop of $\T_L$ an even number of times.
\item Prove that an even subset of $E^*$ has a homology class in
$(\Z/2)^2$ (its parities of crossings with a horizontal and a vertical loop),
that $\mathcal B$ is the class $(0,0)$, and that the even subsets of class
$(a,b)$ are exactly the images $\Gamma(\sigma)$ of configurations with
twisted boundary conditions $\varepsilon=(\varepsilon_1,\varepsilon_2)=((-1)^b,(-1)^a)$
(the coupling across the seam in direction $r$ is $\varepsilon_rK\sigma_i\sigma_j$).
Denote the corresponding partition functions $Z^{\varepsilon}_{\T_L}(K)$.
\item Combining (a) with (b)--(c) on the dual torus (which is again $\T_L$),
prove the Kramers--Wannier identity
\[
 \frac{Z_{\T_L}(K)}{(\sinh2K)^{L^2/2}}=\frac12\sum_{\varepsilon\in\{\pm1\}^2}
 \frac{Z^{\varepsilon}_{\T_L}(K^*)}{(\sinh2K^*)^{L^2/2}} .
\]
\item Prove $\abs{\log Z^\varepsilon_{\T_L}-\log Z_{\T_L}}\le4KL$ and deduce
that $\Phi(K)=\psi(K)-\tfrac12\log\sinh2K$ (with $\psi(K)=\psi(\beta,0)$,
$K=\beta J$) satisfies $\Phi(K)=\Phi(K^*)$.  The map $K\mapsto K^*$ is a
decreasing involution of $(0,\infty)$ with unique fixed point
\[
 K_{\rm sd}=\tfrac12\log(1+\sqrt2)=\tfrac12\operatorname{arsinh}1=0.4407\ldots ,
\]
and if $\psi$ has exactly one point of nonanalyticity in $K\in(0,\infty)$, it
is $K_{\rm sd}$.
\item Prove that on a bipartite graph at $h=0$, $Z(K)=Z(-K)$ and
$\langle\sigma_x\sigma_y\rangle_K=(-1)^{\abs{x-y}_1}\langle\sigma_x\sigma_y\rangle_{-K}$
(reverse the spins of one sublattice), so that the antiferromagnet $J<0$ on
$\Z^2$ has the same pressure and the same critical coupling.
\end{enumerate}
\end{exercise}

\begin{exercise}[Prediction: the square-lattice Ising critical temperature and singularities]
Consider the infinite square-lattice nearest-neighbour Ising
model with $J>0$, $h=0$, and $K=J/(k_BT)$.
For this final exercise assume, as proved by Onsager
\cite{onsager-ising}, the pressure formula
\[
 \psi(K)=\frac12\log(2\sinh2K)
 +\frac1{2\pi}\int_0^\pi
       \operatorname{arcosh}(\cosh2K\,\coth2K-\cos\theta)\dd\theta.
\]
For the plus Gibbs state assume, as proved by Yang
\cite{yang-magnetization}, that
$m^*(K)=(1-\sinh^{-4}2K)^{1/8}$ for
$K>K_{\rm sd}$ and $m^*(K)=0$ for $K\le K_{\rm sd}$,
with $K_{\rm sd}$ from the duality exercise.
\begin{enumerate}[label=(\alph*)]
\item Set $s=\sinh2K$ and use
$\cosh2K\,\coth2K=s+s^{-1}\ge2$ to show that the
pressure is analytic away from $K_{\rm sd}$.
Locate the singular integrand at $\theta=0$ and derive
\[
 K_c=K_{\rm sd}=\tfrac12\log(1+\sqrt2),\qquad
 k_BT_c=2.269185\ldots J.
\]
\item Split the integral at a fixed small $\theta_0$.
Near $(K_c,0)$ use
$\operatorname{arcosh}(\cosh2K\,\coth2K-\cos\theta)
=\sqrt{16(K-K_c)^2+\theta^2}$ to leading order.
Control the differentiated remainder and derive
\[
 \frac c{k_B}=\frac{8K_c^2}{\pi}
       \log\frac1{|1-T/T_c|}+O(1).
\]
Calculate the logarithmic amplitude and its ratio above
and below the transition.
\item Expand the magnetization formula to obtain
$m^*(K)\sim(8\sqrt2)^{1/8}(K-K_c)^{1/8}$ and hence
$\beta_{\rm mag}=1/8$. Use the section Gibbs states to
identify the ordered and unique-state regimes.
\item Under the bipartite spin reversal derive the
corresponding prediction for staggered magnetization in
the square-lattice antiferromagnet. State which fitted
quantities test the exponent and which only set the
nonuniversal energy scale $J$. Separate the model's
strictly two-dimensional, nearest-neighbour, zero-field
assumptions from any material's interlayer couplings or
anisotropy corrections.
\end{enumerate}
\end{exercise}
